\begin{abstract}
		El uso del lenguaje gráfico en la computación cuántica para la representación de algoritmos, aunque intuitivo, no es muy útil para diferentes tareas como la descripción de circuitos cuánticos en entornos de texto, el cálculo de estados cuánticos o la optimización de circuitos cuánticos.
		Mientras que los circuitos clásicos pueden representarse tanto mediante grafos de circuitos como mediante expresiones booleanas, los circuitos cuánticos se han ilustrado hasta ahora predominantemente como grafos de circuitos porque no se ha aceptado hasta el momento ningún lenguaje formal para circuitos cuánticos que permita manipulaciones algebraicas.
		Existe una propuesta para representar circuitos cuánticos de manera conveniente y concisa,
		similar a la forma en que las expresiones booleanas se usan en los circuitos clásicos.
		La notación propuesta permite la descripción consistente y parametrizada de algoritmos cuánticos, así como el manejo fácil de los elementos que la componen para lograr poderosas optimizaciones en el número de puertas de los circuitos.
		Basándonos en esta propuesta, realizamos una extensión de la notación que nos permite manejar más fácilmente un número arbitrariamente grande de índices; para ello introducimos la notación y las operaciones básicas cuando la indicación de los índices está basada en conjuntos.
\end{abstract}
